\subsection{Oscillatory Phasor Contrasts After a Known Restoration}
\label{app:modal-recovery}

The phase and amplitude of an electromechanical oscillation affect the
sequence of fundamental phasors reported by a PMU. Zhou \emph{et al.}
used these variations to estimate oscillation modes from ringdowns
\cite{zhou2010phasorprony}. We tested whether
the shape of this sequence helps distinguish a second switch from the
continuing response to the first. All estimators received the first
restoration time. From then on, each searched for a second event using
only the voltage and current magnitudes at the eight PMUs. The predictor
and the simulation shared the second plant's model and 18 candidate
assets.

For a candidate onset at report $j$, let $Y_-$ and $Y_+$ collect the
30 magnitude reports on either side of the cut. Report $j$ starts
$Y_+$. Both blocks must follow the first restoration. The decision is
available at $j+29$. We projected $Y_+-Y_-$ onto one of three temporal
bases. The first retains the mean. The second also retains sine and
cosine components at \ModalFrequency{} Hz, selected from the low-frequency
eigenvalues of the machine-state Jacobian at equilibrium. The third
uses a slope and curvature, giving it the same dimension as the
sinusoidal basis. We orthonormalized each basis before projection.

Let $z$ collect the projected contrasts after centering and scaling on
48 independent recovery-only records. Let $d_s$ be the projected and
scaled isolated response of a specified asset and switch size, and
$\Sigma$ the estimated covariance of $z$ during recovery. Under the
additive Gaussian approximation, the second switch changes the mean
from zero to $d_s$ and has log likelihood ratio
\[
  \ell_s(z)=d_s^\top\Sigma^{-1}z-
                 \tfrac12 d_s^\top\Sigma^{-1}d_s .
\]
The model averages $\exp(\ell_s)$ over switch sizes within each asset,
assets within each family, and the two families with equal prior weights
at each level. At the first threshold crossing it reports the family
with the greatest marginal mass and the most probable asset in that
family. Each variant's threshold used the same 80 recovery-only
calibration records and 36 long normal records.

For each of 96 new first interventions we simulated two continuations
with the same initial history and instrument noise. One simply
recovered. The other had a second load or line event 1.25--4\,s after
restoration. Each ordered family pair occurred 24 times. A timely
diagnosis had to be published between the onset and clearance of the
second event. Early alarms exhausted the single decision on that record.
The separate normal test measured excursions above threshold across
\ModalQuietHours{} monitored hours.

\begin{table}[!ht]
\centering\small
\caption{Paired continuations after a supplied first restoration.
The first three counts are timely outcomes on 96 second events.
The last counts alarms on 96 recovery-only twins.}
\label{tab:modal-recovery}
\setlength{\tabcolsep}{3pt}
\begin{tabular}{@{}lrrrr@{}}
\toprule
Temporal basis & Detected & Family & Asset & False alarms\\
\midrule
\ModalRows
\bottomrule
\end{tabular}
\end{table}

Oscillatory contrasts detected more second events than the mean alone
at the same calibration budget. More correct family diagnoses came
from load changes, while more lines were classified as loads. The
paired family comparison gave $p=\ModalPMean{}$ against the mean and
$p=\ModalPQuadratic{}$ against slope and curvature. Each estimator
issued zero alarms on the long normal test records (one-sided 95\%
Poisson upper bound \ModalQuietUpper{}/h). Recovery-only records
produced the false alarms in Table~\ref{tab:modal-recovery}.

We then tested whether giving the classifiers time to observe the
transient would improve attribution. A fresh set of 96 paired
continuations used the same design. Every classifier received the
first alarm $b$ of the oscillatory estimator at its original threshold.
At $b+29$ we formed one contrast from the same two blocks for every
candidate. For an onset $d$ intervals before $b$, let $s_a(r)$ denote
the physical response of source and size $a$ at lag $r$, with
$s_a(r)=0$ before the switch. At position $r=0,\ldots,29$ the
predicted contrast is $s_a(r+d)-s_a(r+d-30)$. We evaluated the
likelihood with $d=0$ and also averaged it over $d=0,\ldots,29$.
Uniform weights approximate the unknown alarm delay. The original
alarm and the later diagnosis both had to fall within the true event.

\begin{table}[!ht]
\centering\small
\caption{Classification after a common alarm on 96 new second events.
The first row uses the report at the alarm. All remaining rows use the
same additional 29 reports.}
\label{tab:modal-confirmation}
\setlength{\tabcolsep}{4pt}
\begin{tabular}{@{}llrr@{}}
\toprule
Temporal basis & Assumed onset & Family & Asset\\
\midrule
\ModalConfirmRows
\bottomrule
\end{tabular}
\end{table}

The common alarm detected \ModalTriggerCount{} of the 96 second events
in time and fired on \ModalConfirmNull{} of the 96 recovery-only twins.
The latter exceeded the prespecified 5\% recovery alarm budget on this
cohort. With onset averaging, the sinusoidal estimator identified
\ModalConfirmFamily{} families and \ModalConfirmSite{} exact assets.
The mean matched its family count. The quadratic basis identified
\ModalConfirmQuad{} families (both Holm-adjusted paired $p=1$ versus
the sinusoidal basis). The additional observation time accounted for
most of the family gain. Accounting for an onset before the alarm
helped resolve the asset in all three bases. These simulator-matched
tests supply the true first clearance. An operational monitor must
estimate it.
